\documentclass[11pt]{article}
\usepackage[letterpaper,margin=0.78in]{geometry}
\usepackage{amsmath,amssymb,amsthm,mathtools}
\usepackage{booktabs,enumitem,microtype}
\usepackage[dvipsnames]{xcolor}
\usepackage[colorlinks=true,linkcolor=MidnightBlue,citecolor=MidnightBlue,urlcolor=MidnightBlue]{hyperref}
\setlist{nosep,leftmargin=*}
\newtheorem{theorem}{Theorem}
\newtheorem{lemma}[theorem]{Lemma}
\newtheorem{proposition}[theorem]{Proposition}
\newcommand{\Hh}{\mathcal H}
\newcommand{\C}{\mathcal C}
\newcommand{\R}{\mathbb R}
\title{An explicit upper bound for Voiculescu's constants $\gamma_n$\\
\large A rigorous partial answer to Problem 1 of arXiv:1810.12497}
\author{Autonomous solve-first investigation, lane 12}
\date{11 August 2026}

\begin{document}
\maketitle

\begin{abstract}
Voiculescu asks for quantitative bounds on the universal constants $\gamma_n$
in the formula relating the Lorentz $(n,1)$ modulus of quasicentral
approximation to spectral multiplicity.  A direct dyadic construction gives
the explicit estimate
\[
                 0<\gamma_n\le \frac{2n^n}{12^{n/2}}
                 \qquad(n\ge2).
\]
The proof computes every singular value of the finite-rank commutators.  We
also record the corresponding rectangular scaling and show that increasing
the local tensor-polynomial degree cannot improve this construction.  No
claim of sharpness or new quantitative lower bound is made.
\end{abstract}

\paragraph{The source problem.}
For a compact operator $X$, write
\[
 |X|_n^-=\sum_{r\ge1}s_r(X)r^{-1+1/n},
\]
where $s_1(X)\ge s_2(X)\ge\cdots$ are its singular values.  The modulus
$k_n^-(\tau)$ is the infimum of the limiting values of
$\max_j|[A_m,T_j]|_n^-$ over finite-rank positive contractions
$A_m\uparrow I$.  For a commuting Hermitian $n$-tuple $\tau$, the known
formula is
\[
 (k_n^-(\tau))^n=\gamma_n\int_{\R^n}m(s)\,ds,
 \tag{1}\label{eq:universal}
\]
with $0<\gamma_n<\infty$.  Problem 1 asks for upper and lower bounds for
$\gamma_n$, $n\ge2$ \cite[p.~21]{V}.

\section{The dyadic projection calculation}

Let $\Hh=L^2([0,1]^n)$ and let
$\tau=(M_{x_1},\ldots,M_{x_n})$.  Its absolutely continuous multiplicity is
one on the unit cube and zero elsewhere, so (1) gives
\[
                         \gamma_n=(k_n^-(\tau))^n.       \tag{2}
\]
For $N=2^m$, partition the cube into $N^n$ congruent cubes of side
$h=N^{-1}$, and let $P_N$ be the orthogonal projection onto functions that
are constant on each cube.  These projections are finite rank, nested, and
increase strongly to $I$.

\begin{lemma}[exact cell block]\label{lem:block}
For every coordinate $j$, the operator $[P_N,M_{x_j}]$ has exactly $2N^n$
nonzero singular values, all equal to $1/(\sqrt{12}N)$.
\end{lemma}

\begin{proof}
Fix one cell $Q$ and put
\[
 e=|Q|^{-1/2}\mathbf1_Q,\qquad
 c=|Q|^{-1}\int_Qx_j\,dx,\qquad
 w=(x_j-c)e.
\]
Then $w\perp e$ and $\|w\|_2=h/\sqrt{12}$.  If
$u=w/\|w\|_2$ and $P_Q$ is projection onto $\C e$, then, in the ordered
basis $(e,u)$,
\[
 [P_Q,M_{x_j}]=
 \begin{pmatrix}0&h/\sqrt{12}\\-h/\sqrt{12}&0\end{pmatrix}.
 \tag{3}
\]
Indeed, $M_{x_j}e=ce+w$ and
$\langle M_{x_j}u,e\rangle=\langle u,M_{x_j}e\rangle=h/\sqrt{12}$.
If $v\perp e,u$, then both $P_QM_{x_j}v$ and $M_{x_j}P_Qv$ vanish, so (3)
is the whole nonzero block on $L^2(Q)$.  The blocks for the $N^n$ disjoint
cells form an orthogonal direct sum, proving the claim.
\end{proof}

Consequently, for every $j$,
\[
 |[P_N,M_{x_j}]|_n^-
 =\frac1{\sqrt{12}N}\sum_{r=1}^{2N^n}r^{-1+1/n}.       \tag{4}
\]
The elementary power-sum asymptotic
$\sum_{r=1}^M r^{-1+1/n}\sim nM^{1/n}$ now gives
\[
 \lim_{N\to\infty}|[P_N,M_{x_j}]|_n^-
 =\frac{n2^{1/n}}{\sqrt{12}}.                          \tag{5}
\]

\begin{theorem}[explicit universal upper bound]\label{thm:main}
For every integer $n\ge2$,
\[
                       \boxed{\displaystyle
                       \gamma_n\le\frac{2n^n}{12^{n/2}}.}
\]
\end{theorem}

\begin{proof}
The same $P_N$ works simultaneously for all coordinates, and the quantities
in (4) are equal.  Thus (5) is an admissible sequence in the definition of
$k_n^-(\tau)$, whence
$k_n^-(\tau)\le n2^{1/n}/\sqrt{12}$.  Raise this inequality to the $n$th
power and use (2).
\end{proof}

For orientation, the bound is $\gamma_2\le2/3$,
$\gamma_3\le3\sqrt3/4$, and $\gamma_4\le32/9$.  The construction also
applies when $n=1$, but its estimate $1/\sqrt3$ is weaker than the known
identity $\gamma_1=1/\pi$.

\section{Scaling and a natural optimization test}

\paragraph{Rectangles.}
On $R=\prod_{j=1}^n[0,L_j]$, divide direction $j$ into $N_j$ intervals and
put $M=\prod_jN_j$.  The same proof shows that the $j$th commutator has
$2M$ nonzero singular values, all equal to
$L_j/(\sqrt{12}N_j)$.  Choosing $N_j/L_j$ asymptotically independent of
$j$ yields
\[
 \limsup\max_j|[P,M_{x_j}]|_n^-
 \le \frac{n2^{1/n}}{\sqrt{12}}\,|R|^{1/n}.
\]
Thus the construction has exactly the volume scaling forced by (1), rather
than relying on an accidental symmetry of the unit cube.

\paragraph{Higher local degree.}
One might try to improve the estimate by projecting on each cell onto tensor
polynomials of coordinate degree at most $r$.  The orthonormal shifted
Legendre recurrence shows that only the degree-$r$ boundary couples out of
this subspace.  In each coordinate the commutator therefore has
$2(r+1)^{n-1}N^n$ equal nonzero singular values
\[
 \frac{a_r}{N},\qquad
 a_r=\frac{r+1}{2\sqrt{(2r+1)(2r+3)}}.
\]
Its limiting Lorentz norm is
\[
 C_{n,r}=n\,[2(r+1)^{n-1}]^{1/n}a_r.                   \tag{6}
\]

\begin{proposition}
For $n\ge2$ and every integer $r\ge1$, $C_{n,r}>C_{n,0}$.
Hence piecewise constants are optimal within this tensor-polynomial family.
\end{proposition}

\begin{proof}
After squaring the ratio in (6), the desired inequality is
\[
 \frac{3(r+1)^{,4-2/n}}{(2r+1)(2r+3)}>1.
\]
Since $4-2/n\ge3$ and
$3(r+1)^3-(2r+1)(2r+3)=r(3r^2+5r+2)>0$, the claim follows.
\end{proof}

\paragraph{Scope and remaining obstruction.}
The theorem supplies a fully explicit upper bound, which is one side of the
source problem.  The natural dual approach uses truncated vector Riesz
potentials so that a sum of coordinate commutators has a positive rank-one
part.  Turning that idea into a numerical lower bound requires a sharp,
explicit weak-Schatten estimate for those truncated kernels; that step was
not obtained.  Accordingly, the strict positivity in the abstract is the
previously known part of (1), and this packet does not claim sharpness, a new
quantitative lower bound, or a full solution of the hybrid problem.

\begin{thebibliography}{9}
\bibitem{V}
D.-V. Voiculescu,
\emph{Commutants mod Normed Ideals},
arXiv:1810.12497 (2018; version 2, 2019).
\end{thebibliography}

\end{document}
